\chapter{The finite geometric symmetry group}
\label{ch:ArchonFull600Theory_Symmetry}

This chapter is Milestone 3 of the handoff: O3, O5, Lemma 4 without its
Coxeter presentation (the orders \(7200\) and \(14400\), the kernel
computation and the stabilizer bound from the icosahedral link), and the
first part of DS, the covering of \(SO(4)\) by
\(\operatorname{Sp}(1)\times\operatorname{Sp}(1)\). The presentation
\(K\cong W(H_4)\) is Milestone 4. Source line numbers refer to
\texttt{research/Full\_600cell\_Puzzle\_Theory.tex}.

\section{The binary icosahedral group}

\begin{theorem}[\(V/2\) is a group of order 120]
  \label{thm:binary-icosahedral}
  \lean{ArchonFull600Theory.binaryIcosahedral, ArchonFull600Theory.binaryIcosahedralGroup,
    ArchonFull600Theory.halfVertexSet_closure}
  \leanok
  % SOURCE: Definitions and mathematical sources, lines 280--289
  % SOURCE QUOTE: Identify \(\mathbb R^4\) with the quaternion algebra \(\mathbb H\) using
  % the ordered orthonormal basis \(1,i,j,k\). [...] The binary icosahedral group \(I^*\)
  % is the inverse image of the rotation group of a regular icosahedron under
  % \operatorname{Sp}(1)\longrightarrow SO(3)
  \textit{Source: Definitions (O3); computation: the closure loop of
  \texttt{verify\_regular\_geometry.py}.}
  The set \(V/2\) of \(120\) unit quaternions contains \(1\) and \(-1\) and
  is closed under the quaternion product and under conjugation. With the
  group structure of \(\operatorname{Sp}(1)\) it is therefore a subgroup
  \(I^*\) with \(120\) elements, which is how \(I^*\) is defined here.
  That this subgroup is the inverse image of the icosahedral rotation
  group is E1, which stays a hypothesis
  (Chapter~\ref{ch:ArchonFull600Theory_Geometry}).
\end{theorem}
\begin{proof}
  \leanok
  Exact products in \(\mathbb Z[\phi]\) show that left multiplication by
  the vertices \(8\) and \(30\) permutes \(V/2\), with the two permutation
  tables checked by the kernel. A breadth-first tree of words in these two
  vertices, with a parent table and strictly decreasing word lengths,
  writes every element of \(V/2\) as such a word; induction on the word
  length shows that left multiplication by any element of \(V/2\) permutes
  \(V/2\). A third table shows that conjugation permutes \(V/2\).
\end{proof}

\section{The maps \texorpdfstring{\(L_{p,q}\)}{L(p,q)} and \texorpdfstring{\(J\)}{J}}

\begin{definition}[\(L_{p,q}\) and \(J\)]
  \label{def:lpq}
  \lean{ArchonFull600Theory.lpq, ArchonFull600Theory.lpqHom, ArchonFull600Theory.conjIso,
    ArchonFull600Theory.conjIso_lpq_conjIso}
  \leanok
  % SOURCE: Lemma 4, lines 524--527, 547, 561--567 and 581--585
  % SOURCE QUOTE: L_{p,q}(x)=pxq^{-1}. [...] Quaternion norm multiplicativity makes
  % \(L_{p,q}\) orthogonal. [...] J(x)=\bar x [...] J L_{p,q}J=L_{q,p},
  \textit{Source: Lemma 4 (O5).}
  For unit quaternions \(p,q\), the map \(L_{p,q}(x)=pxq^{-1}=px\bar q\) is
  an orthogonal map of \(\mathbb H\), and \((p,q)\mapsto L_{p,q}\) is a
  homomorphism from \(\operatorname{Sp}(1)\times\operatorname{Sp}(1)\).
  Quaternion conjugation \(J(x)=\bar x\) is orthogonal, and
  \(JL_{p,q}J=L_{q,p}\).
\end{definition}

\begin{lemma}[The kernel]
  \label{lem:lpq-kernel}
  \lean{ArchonFull600Theory.lpq_eq_one_iff, ArchonFull600Theory.ker_lpqHom}
  \leanok
  \uses{def:lpq}
  % SOURCE: Lemma 4, lines 553--558
  % SOURCE QUOTE: If \(L_{p,q}\) is the identity, evaluating at \(1\) gives \(p=q\). This
  % quaternion must then commute with every quaternion. Commuting with \(i\)
  % and \(j\) forces it to be real, so it is \(1\) or \(-1\). The kernel is
  % precisely \{(1,1),(-1,-1)\}.
  \textit{Source: Lemma 4, kernel calculation.}
  \(L_{p,q}\) is the identity exactly when \((p,q)=(1,1)\) or
  \((p,q)=(-1,-1)\).
\end{lemma}
\begin{proof}
  \leanok
  Evaluating at \(1\) gives \(p=q\), so \(p\) commutes with every
  quaternion. The \(j\)- and \(k\)-components of \(pi=ip\) and the
  \(k\)-component of \(pj=jp\) make the imaginary part zero, and
  \(|p|=1\) gives \(p=\pm1\).
\end{proof}

\begin{lemma}[Determinants]
  \label{lem:lpq-det}
  \lean{ArchonFull600Theory.orthDet_lpq, ArchonFull600Theory.orthDet_conjIso,
    ArchonFull600Theory.det_mulLeft}
  \leanok
  \uses{def:lpq}
  % SOURCE: Lemma 4, lines 547--550 and 563--567
  % SOURCE QUOTE: Its determinant is positive: the same construction on the connected space
  % \(\operatorname{Sp}(1)\times\operatorname{Sp}(1)\) depends continuously on \((p,q)\)
  % and includes the identity. [...] Its real matrix is
  % \operatorname{diag}(1,-1,-1,-1), so \(\det J=-1\).
  \textit{Source: Lemma 4.}
  \(\det L_{p,q}=1\) for unit quaternions \(p,q\), and
  \(\det J=-1\).
\end{lemma}
\begin{proof}
  \leanok
  In the basis \(1,i,j,k\), left multiplication by \(a\) has determinant
  \(|a|^4\), and \(J\) has matrix \(\operatorname{diag}(1,-1,-1,-1)\).
  Since \(L_{p,q}(x)=p\,\overline{q\,\bar x}\), its determinant is
  \(|p|^4\cdot(-1)\cdot|q|^4\cdot(-1)=1\). This replaces the source's
  continuity argument by a direct computation.
\end{proof}

\section{The covering of \texorpdfstring{\(SO(4)\)}{SO(4)}}

\begin{theorem}[Orthogonal maps of \(\mathbb H\)]
  \label{thm:orthogonal-lpq}
  \lean{ArchonFull600Theory.orthogonal_eq_lpq, ArchonFull600Theory.orthDet_eq_one_or,
    ArchonFull600Theory.reflection_mem_quaternionForms}
  \leanok
  \uses{def:lpq, lem:lpq-det}
  % SOURCE: Distinct symmetry and motion groups, lines 1350--1355; handoff E6
  % SOURCE QUOTE: To see the continuous covering directly, take \(A\in SO(4)\) and put
  % \(p=A(1)\). [...] Every such three-dimensional rotation is \(x\mapsto qxq^{-1}\)
  % for a unit quaternion \(q\), so \(A\) has a unit-quaternion-pair representation.
  \textit{Source: Distinct groups (DS), with E6 proved rather than assumed.}
  Every orthogonal map of \(\mathbb H\) is \(L_{p,q}\) or \(L_{p,q}J\) for
  unit quaternions \(p,q\); in particular its determinant is \(1\) or
  \(-1\).
\end{theorem}
\begin{proof}
  \leanok
  The reflection in the hyperplane orthogonal to \(v\neq0\) is
  \(x\mapsto-v\bar xv/|v|^2=L_{u,-\bar u}J(x)\) with \(u=v/|v|\), from the
  identity \(v\bar xv=2\langle v,x\rangle v-|v|^2x\). The maps \(L_{p,q}\)
  and \(L_{p,q}J\) are closed under composition, since
  \(JL_{p,q}=L_{q,p}J\) and \(J^2=1\). Every orthogonal map of a
  finite-dimensional real inner product space is a product of reflections
  (the Cartan--Dieudonn\'e theorem, in Mathlib). This replaces the source's
  route through \(SO(3)\).
\end{proof}

\begin{theorem}[DS, part 1: the double cover]
  \label{thm:double-cover}
  \lean{ArchonFull600Theory.doubleCoverSO4, ArchonFull600Theory.specialOrthogonal}
  \leanok
  \uses{lem:lpq-kernel, lem:lpq-det, thm:orthogonal-lpq}
  % SOURCE: Distinct symmetry and motion groups, lines 1339--1340 and 1350--1355
  % SOURCE QUOTE: \(\operatorname{Sp}(1)\times\operatorname{Sp}(1)\) & Double cover of
  % \(SO(4)\), with diagonal \(\{\pm(1,1)\}\) kernel [...] The kernel calculation is the
  % same as in Lemma 4.
  \textit{Source: Distinct groups (DS), covering part.}
  The homomorphism \((p,q)\mapsto L_{p,q}\) maps
  \(\operatorname{Sp}(1)\times\operatorname{Sp}(1)\) onto \(SO(4)\), the
  orthogonal maps of \(\mathbb H\) of determinant \(1\), and its kernel is
  \(\{(1,1),(-1,-1)\}\).
\end{theorem}
\begin{proof}
  \leanok
  The image lies in \(SO(4)\) by the determinant lemma. An element of
  \(SO(4)\) is \(L_{p,q}\) or \(L_{p,q}J\), and the second form has
  determinant \(-1\). The kernel is the kernel lemma.
\end{proof}

\section{The groups \texorpdfstring{\(K\)}{K} and \texorpdfstring{\(K^+\)}{K+}}

\begin{definition}[\(K\) and \(K^+\)]
  \label{def:symmetry-group}
  \lean{ArchonFull600Theory.symmetryGroup, ArchonFull600Theory.rotationGroup,
    ArchonFull600Theory.mem_symmetryGroup_iff, ArchonFull600Theory.isometry_mem_symmetryGroup}
  \leanok
  % SOURCE: Overview table, lines 236--239
  % SOURCE QUOTE: \(K^+\) & 7,200 elements & Orientation-preserving rigid symmetries of
  % the whole polytope \\ \(K=W(H_4)\) & 14,400 elements & Full rigid symmetry group,
  % including reflections
  \textit{Source: overview table (O5).}
  \(K\) is the group of orthogonal maps of \(\mathbb H\) that map
  \(P=\operatorname{conv}V\) onto itself, and \(K^+=K\cap SO(4)\). An
  orthogonal map lies in \(K\) exactly when it maps \(V\) onto itself. Every
  isometry of \(\mathbb H\) that maps \(P\) onto itself is an element
  of \(K\) (in particular it fixes \(0\)), so \(K\) is the full rigid
  symmetry group.
\end{definition}
\begin{proof}
  \leanok
  \(V\) is the set of extreme points of \(P\), and linear bijections
  preserve extreme points and convex hulls. An isometry \(f\) with
  \(f(P)=P\) maps the closed ball of radius \(2\) about \(0\), which
  contains \(P\), to one about \(c=f(0)\), so \(|\pm v-c|\le2\) for the
  vertex \(v=2\); then \(8+2|c|^2=|v-c|^2+|v+c|^2\le8\) gives \(c=0\), and
  the Mazur--Ulam theorem makes \(f\) linear.
\end{proof}

\begin{lemma}[Stabilizer bound]
  \label{lem:stabilizer-bound}
  \lean{ArchonFull600Theory.card_stabilizer_le, ArchonFull600Theory.stabilizer_injective,
    ArchonFull600Theory.card_linkTriangles, ArchonFull600Theory.eq_of_eqOn_clique}
  \leanok
  \uses{def:symmetry-group}
  % SOURCE: Lemma 4, lines 569--578
  % SOURCE QUOTE: An upper bound is still necessary. Fix a vertex \(v\). Its stabilizer
  % acts faithfully on the icosahedral vertex figure: the vertex direction
  % and the three-dimensional span of that figure together span
  % \(\mathbb R^4\). An orthogonal symmetry of an icosahedron is determined
  % by the image of one ordered triangular face. There are at most
  % 20\cdot6=120 such images.
  \textit{Source: Lemma 4, upper bound.}
  The stabilizer in \(K\) of vertex \(0\) has at most \(20\cdot6=120\)
  elements, and \(|K|\le120\cdot120=14400\).
\end{lemma}
\begin{proof}
  \leanok
  \(K\) permutes the vertices and preserves the edge graph. A symmetry
  fixing vertex \(0\) maps the facet \(\{0,30,31,32\}\) to vertex \(0\)
  and an ordered triangle of pairwise adjacent neighbours of vertex
  \(0\). The four vertices of a facet are linearly independent, so they
  are a basis of \(\mathbb H\) and the symmetry is determined by these
  images. The twelve neighbours of vertex \(0\) have \(120\) ordered
  triangles, a kernel count. Orbit--stabilizer with an orbit of at most
  \(120\) vertices gives the bound on \(|K|\), and also that \(K\) is
  finite.
\end{proof}

\begin{theorem}[Lemma 4 without the presentation]
  \label{thm:finite-symmetry-group}
  \lean{ArchonFull600Theory.finiteSymmetryGroup, ArchonFull600Theory.rotationGroupEquiv,
    ArchonFull600Theory.card_rotationGroup, ArchonFull600Theory.card_symmetryGroup,
    ArchonFull600Theory.mem_symmetryGroup_iff_lpq, ArchonFull600Theory.ker_icosianRotation}
  \leanok
  \uses{thm:binary-icosahedral, lem:lpq-kernel, lem:lpq-det, lem:stabilizer-bound}
  % SOURCE: Lemma 4, lines 524--531 and 550--585
  % SOURCE QUOTE: K^+\cong (I^*\times I^*)/\langle(-1,-1)\rangle,
  % \qquad |K^+|=7200, \end{equation} and quaternion conjugation extends this to \(K\) of
  % order \(14400\). [...] The exhibited maps attain this upper bound, proving both
  % group orders and the description of \(K^+\).
  \textit{Source: Lemma 4.}
  \(K^+\cong(I^*\times I^*)/\langle(-1,-1)\rangle\) through
  \((p,q)\mapsto L_{p,q}\), and \(|K^+|=7200\). \(J\) lies in \(K\) but not
  in \(K^+\), and \(|K|=14400\). Every element of \(K\) is \(L_{p,q}\) or
  \(L_{p,q}J\) with \(p,q\in I^*\).
\end{theorem}
\begin{proof}
  \leanok
  Closure of \(I^*\) under products and conjugation shows that
  \(L_{p,q}\) and \(J\) map \(V\) to itself. On \(I^*\times I^*\), which has
  \(14400\) elements, the kernel of \((p,q)\mapsto L_{p,q}\) is
  \(\{(1,1),(-1,-1)\}\), so its image in \(K^+\) has \(7200\) elements. The
  maps \(A\mapsto A\) and \(A\mapsto AJ\) embed two disjoint copies of
  \(K^+\) in \(K\), so \(2|K^+|\le|K|\le14400\). Hence the image is all of
  \(K^+\), \(|K^+|=7200\) and \(|K|=14400\). An element of \(K\) of
  determinant \(-1\) is \(AJ\) with \(A=(AJ)J\in K^+\).
\end{proof}
